\section{Background: From RLHF to RLVR and the GRPO Family}\label{sec:background}

We read the algorithms behind nearly all LLM-agent training through one question: \emph{at what granularity does each method assign credit?} The field moved from a learned per-token signal (a critic) to a per-response scalar shared by every token, and the 2025--2026 refinements left this property intact.

\paragraph{RLHF with a critic}
Classic RLHF \citep{christiano2017deep,ouyang2022training} fits a reward model (RM) to human preferences and optimizes the policy against it with PPO \citep{schulman2017proximal} plus a per-token KL penalty. The RM emits one scalar per response; the PPO critic spreads it over tokens through generalized advantage estimation (GAE; \citealp{schulman2015high}; Eq.~\ref{eq:gae} in Section~\ref{sec:problem}). This is genuine token-level credit, but it keeps four models in memory (policy, reference, RM, critic) and requires training the critic. DPO \citep{rafailov2023direct} removes the RM and the sampling loop at the price of being offline; a careful comparison nonetheless found well-tuned PPO ahead of DPO, especially on code \citep{xu2024dpo}\unverified.

\paragraph{Why critic-free methods took over}
For long chain-of-thought (CoT) reasoning the critic turned out to be not just expensive but \emph{wrong}. VC-PPO \citep{yuan2025whats} identifies \emph{value-initialization bias} and \emph{reward-signal decay}: with $\lambda<1$ a terminal reward is attenuated geometrically over thousands of tokens and barely reaches the early ones. VAPO \citep{yue2025vapo} reports vanilla PPO at roughly 5 points on AIME24 versus 60.4 for VAPO (Qwen2.5-32B base). Value-based methods can thus be repaired, but the mainstream took the other branch: drop the critic and accept an unbiased, high-variance, \emph{uniform} sequence-level signal. RLOO \citep{ahmadian2024back} argued that, from a strong SFT policy, RLHF is effectively a bandit whose single action is the whole response\unverified{} (Table~\ref{tab:grpo-variants}, row a).

\paragraph{RLVR and GRPO}
RL with verifiable rewards (RLVR; popularized by \citealp{lambert2024tulu}) replaces the RM with a programmatic checker, such as exact match of a math answer or a unit-test suite, returning $r\in\{0,1\}$. GRPO \citep{shao2024deepseekmath}, the optimizer behind DeepSeek-R1 \citep{deepseekai2025deepseek}, samples a group of $G$ responses $o_1,\dots,o_G$ for a prompt $q$ and maximizes
\begin{equation}\label{eq:grpo}
\mathcal{J}(\theta)=\mathbb{E}\Bigg[\frac{1}{G}\sum_{i=1}^{G}\frac{1}{|o_i|}\sum_{t=1}^{|o_i|}\Big(\min\big(\rho_{i,t}\hat A_{i,t},\,\mathrm{clip}(\rho_{i,t},1-\epsilon,1+\epsilon)\,\hat A_{i,t}\big)-\beta\,\mathbb{D}_{\mathrm{KL}}\big[\pi_\theta\,\|\,\pi_{\mathrm{ref}}\big]\Big)\Bigg],
\end{equation}
where $\rho_{i,t}$ is the token-level importance ratio between $\pi_\theta$ and $\pi_{\theta_{\mathrm{old}}}$ and, under outcome supervision,
\begin{equation}\label{eq:grpo-adv}
\hat A_{i,t}\;=\;\frac{r_i-\mathrm{mean}(r_1,\dots,r_G)}{\mathrm{std}(r_1,\dots,r_G)}\qquad\text{for every } t=1,\dots,|o_i|.
\end{equation}
The group mean replaces the critic. The property that matters here is visible in the subscript: $\hat A_{i,t}$ depends on $(q,i)$ but not on $t$. Every token of $o_i$, whether it belongs to the decisive step, a wasted tool call or formatting boilerplate, is pushed by the same amount. Credit is disentangled only statistically, across many samples in which the decisive step varies independently of the rest. (DeepSeekMath also defines a process-supervision variant that sums normalized step rewards from $t$ onward\unverified; R1-style training does not use it.) Outcome rewards add a second, structural loss: a group whose $G$ rewards are all equal has zero advantage and yields no gradient (Eq.~\ref{eq:zerovar}).

\paragraph{The 2025--2026 variants}
Table~\ref{tab:grpo-variants} groups the main variants by the part of the objective they touch: (a) the advantage estimator, (b) how per-token losses are aggregated, and (c) the trust region or off-policy correction. Dr.~GRPO \citep{liu2025understanding} shows that both normalizers in Eqs.~\ref{eq:grpo}--\ref{eq:grpo-adv} are biased: $1/|o_i|$ inflates the length of incorrect responses, and $1/\mathrm{std}$ up-weights nearly all-correct or all-wrong prompts. DAPO \citep{yu2025dapo} bundles four fixes and reaches 50 on AIME24 with Qwen2.5-32B in half the steps of the R1-Zero-Qwen-32B recipe. GSPO \citep{zheng2025group} makes even the importance ratio sequence-level, arguing that the optimization unit should match the unit that receives reward. The right-hand column is the point: only the value-based revivals (VC-PPO, VAPO) change \emph{which} tokens receive credit, and they do so by reinstating a critic.

\begin{table}[t]
\centering
\small
\caption{Main GRPO-family variants (2025--2026). ``Same $\hat A_i$'' means every token of response $i$ receives one shared advantage.}
\label{tab:grpo-variants}
\begin{tabularx}{\textwidth}{@{}l X l@{}}
\toprule
Method & Mechanism (one line) & Who gets credit \\
\midrule
\multicolumn{3}{@{}l}{\emph{(a) Advantage estimator}}\\
RLOO \citep{ahmadian2024back} & Baseline for sample $i$ = mean reward of the other $k-1$ samples & same $\hat A_i$ \\
REINFORCE++ \citep{hu2025reinforce} & No critic; normalize advantages over the global batch, not per group & same $\hat A_i$ \\
Dr.~GRPO \citep{liu2025understanding} & Drop the $1/\mathrm{std}$ (difficulty) normalizer & same $\hat A_i$ \\
DAPO dyn.\ sampling \citep{yu2025dapo} & Discard all-0/all-1 groups; oversample until the batch is full & same $\hat A_i$ \\
VC-PPO \citep{yuan2025whats} & Value pretraining; decoupled GAE ($\lambda_{\text{critic}}=1$, separate actor $\lambda$) & per token (critic) \\
VAPO \citep{yue2025vapo} & VC-PPO plus length-adaptive GAE and DAPO tricks & per token (critic) \\
\midrule
\multicolumn{3}{@{}l}{\emph{(b) Loss aggregation / weighting}}\\
DAPO token-level loss & Average over all tokens in the batch, so long responses are not down-weighted & same $\hat A_i$, reweighted \\
Dr.~GRPO & Remove the $1/|o_i|$ length normalizer & same $\hat A_i$, reweighted \\
ScaleRL \citep{khatri2025art} & Prompt-level loss averaging; batch-level advantage normalization & same $\hat A_i$ \\
\midrule
\multicolumn{3}{@{}l}{\emph{(c) Trust region / off-policy correction}}\\
DAPO Clip-Higher & Decoupled $\epsilon_{\text{low}}<\epsilon_{\text{high}}$ to counter entropy collapse & same $\hat A_i$ \\
CISPO \citep{minimax2025minimax} & Clip the stop-gradient IS weight, not the update; every token keeps a gradient & same $\hat A_i$ \\
GSPO \citep{zheng2025group} & Length-normalized sequence-level ratio and clipping & same $\hat A_i$ \\
SAPO \citep{gao2025soft} & Temperature-controlled soft gate instead of a hard clip; asymmetric for negatives & same $\hat A_i$ \\
\bottomrule
\end{tabularx}
\end{table}

\paragraph{What scaling studies say}
ScaleRL \citep{khatri2025art} spent more than 400k GPU-hours fitting sigmoidal compute--performance curves and found that most knobs (normalization, aggregation, curriculum) shift efficiency rather than the asymptote; only the loss type mattered (CISPO and GSPO above DAPO). Refinements of \emph{how} a uniform signal is optimized change how fast one reaches the ceiling, not the ceiling; the information content of the signal, which credit assignment changes, is the lever these studies did not pull.

\paragraph{Critiques of RLVR}
Three findings suggest that ``reward went up'' is weak evidence that credit went to the right behavior. \citet{yue2025does} find that RLVR models win pass@$k$ at small $k$ but lose to their base models at large $k$, \ie RLVR mostly reweights solutions already in the base model's support (ProRL \citep{liu2025prorl}\unverified{} argues that prolonged, KL-controlled training can push past this boundary). \citet{cui2025entropy} show that, without intervention, entropy is traded for reward along $R=-a\,e^{H}+b$, driven by the covariance between a token's log-probability and its advantage; relatedly, updating only the roughly 20\% of high-entropy ``forking'' tokens is reported to match full updates (\citealp{wang2025beyond}\unverified; Section~\ref{sec:tax-hindsight}). Finally, \citet{shao2025spurious} obtain +21.4 points on MATH-500 for Qwen2.5-Math-7B with \emph{random} rewards (versus +29.1 with ground truth), an effect that does not transfer to Llama3 or OLMo2 and disappears when clipping is disabled. Two lessons follow: a minority of tokens carries most of the useful signal, yet uniform broadcast spends gradient on all of them equally; and any credit-assignment claim evaluated on one model family needs random-reward and cross-family controls.

\paragraph{Bridge to agents}
Everything above treats the response as one action. In a multi-turn agent the ``response'' interleaves model turns and tool observations, the decisive event is often a single tool call or the reading of a single tool result, and Eq.~\ref{eq:grpo-adv} assigns the final outcome to every turn alike. Section~\ref{sec:agentic} describes this setting.
